
\chapter{Blocos de Código e Terminal}

Este capítulo apresenta exemplos de código em dois regimes editoriais: blocos escuros, adequados para simular uma sessão de terminal ou uma janela de execução, e blocos claros, adequados para leitura de código-fonte em materiais didáticos. O template usa \texttt{listings} com \texttt{tcolorbox}, evitando \texttt{minted} para preservar portabilidade sem \texttt{shell-escape}.

\section{Bloco pequeno de código em terminal}

\begin{VSTerminalCode}[title={Terminal pequeno}]
def normalize(value, scale):
    if scale == 0:
        return 0.0
    return value / scale
\end{VSTerminalCode}

\clearpage
\section{Bloco médio de código em terminal}

\begin{VSTerminalCode}[title={Terminal médio}]
def summarize_jobs(jobs):
    total_gpus = 0
    total_nodes = 0
    failed_jobs = []

    for job in jobs:
        total_gpus += job.get("gpus", 0)
        total_nodes += job.get("nodes", 0)
        if job.get("state") == "FAILED":
            failed_jobs.append(job.get("id"))

    return {
        "jobs": len(jobs),
        "gpus": total_gpus,
        "nodes": total_nodes,
        "failed": failed_jobs,
    }
\end{VSTerminalCode}

\clearpage
\section{Bloco grande de código em terminal}

\begin{VSTerminalCode}[title={Terminal grande}]
class ClusterPolicy:
    def __init__(self, max_gpus_per_user, reserved_nodes):
        self.max_gpus_per_user = max_gpus_per_user
        self.reserved_nodes = set(reserved_nodes)

    def can_schedule(self, request, usage):
        user = request["user"]
        required_gpus = request.get("gpus", 0)
        candidate_nodes = request.get("nodes", [])

        if usage.get(user, 0) + required_gpus > self.max_gpus_per_user:
            return False, "gpu quota exceeded"

        for node in candidate_nodes:
            if node in self.reserved_nodes and not request.get("admin", False):
                return False, "reserved node requested"

        if not self._has_required_topology(candidate_nodes):
            return False, "topology constraint not satisfied"

        return True, "accepted"

    def _has_required_topology(self, nodes):
        if len(nodes) <= 1:
            return True
        racks = {node.split("-")[0] for node in nodes}
        return len(racks) <= 2

policy = ClusterPolicy(max_gpus_per_user=32, reserved_nodes=["rack9-node01"])
request = {
    "user": "researcher",
    "gpus": 8,
    "nodes": ["rack1-node03", "rack1-node04"],
    "admin": False,
}
usage = {"researcher": 16}
accepted, reason = policy.can_schedule(request, usage)
print(accepted, reason)
\end{VSTerminalCode}

\clearpage
\section{Bloco pequeno de código claro}

\begin{VSCodeBlock}[title={Código pequeno}]
def bandwidth(bytes_written, seconds):
    return bytes_written / seconds
\end{VSCodeBlock}

\section{Bloco médio de código claro}

\begin{VSCodeBlock}[title={Código médio}]
def moving_average(samples, window):
    if window <= 0:
        raise ValueError("window must be positive")

    result = []
    for index in range(len(samples)):
        start = max(0, index - window + 1)
        segment = samples[start:index + 1]
        result.append(sum(segment) / len(segment))

    return result
\end{VSCodeBlock}

\clearpage
\section{Bloco grande de código claro}

\begin{VSCodeBlock}[title={Código grande}]
from dataclasses import dataclass

@dataclass
class Measurement:
    timestamp: float
    node: str
    metric: str
    value: float

class TelemetryBuffer:
    def __init__(self, limit):
        self.limit = limit
        self.samples = []

    def append(self, measurement):
        self.samples.append(measurement)
        if len(self.samples) > self.limit:
            self.samples = self.samples[-self.limit:]

    def select(self, metric, node=None):
        selected = []
        for sample in self.samples:
            if sample.metric != metric:
                continue
            if node is not None and sample.node != node:
                continue
            selected.append(sample)
        return selected

    def aggregate_mean(self, metric, node=None):
        selected = self.select(metric, node=node)
        if not selected:
            return None
        return sum(sample.value for sample in selected) / len(selected)

buffer = TelemetryBuffer(limit=10000)
buffer.append(Measurement(0.0, "node001", "gpu_util", 71.0))
buffer.append(Measurement(1.0, "node001", "gpu_util", 78.0))
buffer.append(Measurement(2.0, "node001", "gpu_util", 82.0))
print(buffer.aggregate_mean("gpu_util", node="node001"))
\end{VSCodeBlock}
